% Document Type: LaTeX
% Master File: notation-chk.tex
% definition d'une notation: typiquement
% \notation{notre abbreviation} : string
%          {nombre d'arguments} : 0..9
%          {definition \Latex}  : TeX
%          {Commentaire}        : blabla

 % pour les notations du livre il faut que l'on obtienne la forme suivante
 % | notation | commentaire | usage typique | reference dans le livre |
 % par exemple
 % | =        | equality    | x = y         | def 3 p 345             |

% La macros de definition:
% qui doit etre utilisee dans un tabular de presentation
% \newcommand{\notation}[3]{\verb@#1@ & \verb@#2@ & #3 \\}

\newcommand{\notation}[4]{\newcommand{#1}[#2]{#3}}

%\newcommand{\notation}[4]{{\newcommand{#1}[#2]{#3}}#3 & #4 \\}

%\begin{tabular}{ll}

%%%%%%%%%%%%%%%% Les notations de symboles d'ensemble

\notation{\X}{0}{{\cal X}}{Set of variable symbols}
%\notation{\S}{0}{{\cal S}}{Set of sort symbols}
\notation{\F}{0}{{\cal F}}{Set of function symbols}
\notation{\Fn}{0}{{\cal F}_n}{Set of function symbols of arity $n$}
%\notation{\P}{0}{{\cal P}}{Set of predicate symbols}
\notation{\A}{0}{{\cal A}}{An algebra}
\notation{\Terms}{2}
         {{\calT(#1,#2)}}
         {Set of terms build on variables $X$ and symbols $F$. 
          Typical use $\Terms{F}{X}$}
\notation{\TFX}{0}
         {{\cal T(F,X)}}
         {Set of terms build on variables $\X$ and symbols $\F$}
\notation{\TFV}{0}
         {{\cal T(F,V)}}
         {Set of terms build on variables $\V$ and symbols $\F$}
\notation{\TSV}{0}
         {{\cal T}({\bf \Sigma,\calV})}
         {Set of terms build on variables $\calV$ and signature $\Sigma$}
\notation{\TSX}{0}
         {{\cal T} (\Sigma, {\cal X})}
         {Set of terms build on variables $\bf X$ and signature $\Sigma$}
\notation{\TFXE}{0}
         {{\cal T(F,X)}/E}
         {Quotient of $\TFX$ by the congruence generated by $E$}
\notation{\TSXE}{0}
         {{\cal T}(\Sigma,{\cal X})/E}
         {Quotient of $\TSX$ by the congruence generated by $E$}
\notation{\TF}{0}
         {{\cal T(F)}}
         {Set of ground terms}
\notation{\TS}{0}
         {{\cal T}(\Sigma)}
         {Set of many-sorted ground terms}
\notation{\TFE}{0}
         {{\cal T(F)}/E}
         {Quotient of $\TF$ by the congruence generated by $E$}
\notation{\TSE}{0}
         {{\cal T}(\Sigma)/E}
         {Quotient of $\TS$ by the congruence generated by $E$}
\notation{\TSFX}{0}
         {{\cal T(S,F,X)}}
         {Set of terms build on variables $\X$, sorts $\S$ and symbols $\F$}
\notation{\TSF}{0}
         {{\cal T(S,F)}}
         {Set of terms build on sorts $\S$ and symbols $\F$}
\notation{\FXA}{0}
         {\langle \F,\X,\calA\rangle}
         {Notation for a \FXA-unification problem}
\notation{\Subst}{0}
         {Subst}
         {Set of substitutions of the term algebra}
\notation{\SubstFX}{0}
         {Subst^{\F\X}}
         {Set of substitutions of the term algebra over $\F$ and $\X$}
\notation{\Perm}{0}
         {Perm}
         {Set of permutations, i.e. one to one substitutions of the 
          term algebra}
\notation{\Red}{1}
         {{\cal R}edex(#1)}
         {The set of redexes of a term $u$. Typical use: $\Red{t}$}
\notation{\N}{0}{\mathop{\bf N}}{Set of natural numbers: A eviter,
                prendre plutot la notation Nat}
\notation{\Nat}{0}{{\bf N}}{Set of natural numbers}
\notation{\Reals}{0}{\mathop{\bf R}}{Set of reals}
\notation{\Sol}{2}{{\calS}ol_{#1}(#2)}
         {The set of solutions of a constraint $c$ on the domain $D$.
          Typical use: $\Sol{D}{c}$}
\notation{\SolEqGen}{2}{\Sigma^{\cal U}_{#1}(#2)}
         {Unifiers set of general equation. Typical use:
         $SolEqGen{c}{f,g,A,V,W}$ in the case of collapse axioms,
         $SolEqGen{nc}{f,g,A,V,W}$ in the case of non collapse axioms.}
\notation{\SolMEqGen}{1}{\Sigma^{\cal M}(#1)}
         {Matchers set of general equation. Typical use:
         $SolMEqGen{f,A,V}$}
\notation{\UnifSet}{2}{\calU_{#1}({#2})}
         {Set of $\#1$-unifier of $\#2$. Typical use: $\UnifSet{E}{S}$.}

%%%%%%%%%%%%%%%% symboles d'e'quations et de syste`mes

\notation{\eqq}{0}{=^?}{equation symbol}
\notation{\eqqNT}{0}{=^{{\scriptscriptstyle \neq\topoc} ?}}
         {Restricted equation symbol}
\notation{\eqqO}{0}{=^?_{\emptyset}}{equation symbol in empty theory}
\notation{\eqqE}{0}{=^?_E}{equation symbol modulo $E$}
\notation{\eqqA}{0}{=^?_A}{equation symbol modulo $A$}
\notation{\eqqK}{0}{=^?_K}{equation symbol in the interpretation $\calK$}
\notation{\eqqAC}{0}{=^?_{AC}}{equation symbol modulo $AC$}
\notation{\eqqC}{0}{=^?_{{\sf C}}}{equation symbol modulo $C$}
\notation{\eqqR}{0}{=^?_{{R}}}{equation symbol modulo $R$}
\notation{\meqq}{0}{\ll^?}{match-equation symbol}
\notation{\meqqAC}{0}{\ll^?_{AC}}{AC-match-equation symbol}
\notation{\diseqq}{0}{\neq^?}{disequation symbol}
\notation{\diseqqO}{0}{\neq^?_{\emptyset}}{disequation symbol in empty
          theory}
\notation{\diseqqK}{0}{\neq^?_K}{disequation symbol in the interpretation
          $\calK$}
\notation{\diseqqA}{0}{\neq^?_A}{disequation symbol modulo $A$}
%\notation{\ineqq}{0}{\prec^?}{inequation symbol}
\notation{\succeqq}{0}{\succ^?}{greater than inequation symbol}
\notation{\ineqq}{0}{>^?}{greater than inequation symbol}
\notation{\ineqqO}{0}{>^?_{\emptyset}}{inequation symbol in empty
          theory}
\notation{\ineqqA}{0}{>^?_{A}}{inequation symbol modulo $A$}
\notation{\ineqqK}{0}{>^?_{K}}{inequation symbol in the interpretation
$\calK$}
\notation{\memeqq}{0}{\in^?}{membershipequation symbol}
\notation{\ww}{0}{\:\wedge\:}{connector between equations in systems}
\notation{\vv}{0}{\:\vee\:}{connector between systems in disjunctions}
\notation{\fF}{0}{{\mathbb F}}{the insatisfiable system}
\notation{\tT}{0}{{\mathbb T}}{the always satisfiable system}

%%%%%%%%%%%%%%%% les symboles de selection

\notation{\Var}{1}
         {{\cal V}ar(#1)}
         {Set of variables of a term or an equation or a system}
\notation{\Bvar}{1}
         {{\cal BV}ar(#1)}
         {Bound variables of a term or an equation or a system}
\notation{\Dom}{1}
         {{\cal D}om(#1)}
         {The domain of a substitution or of a term}
\notation{\Ran}{1}
         {{\cal R}an(#1)}
         {The range of a substitution}
\notation{\Grd}{1}
         {{\cal G}rd(#1)}
         {The non-variable positions of a term}
\notation{\restrict}{2}
         {{#1}_{|#2}}
         {Restriction of a substitution to a set of variables}
\notation{\size}{1}
         {|#1|}
         {Size of a term or of an equation}
\notation{\card}{1}
         {|#1|}
         {Cardinality of a set or of a multiset}
\notation{\topoc}{0}
%        {\mbox{\footnotesize $\Lambda$}}
         {{\scriptstyle \Lambda}}
         {The name of the top occurrence of a term}
\notation{\topocc}{0}
         {\topoc}
         {The name of the top occurrence of a term}
\notation{\atoc}{1}
         {_{|#1}}
         {Selection of an occurrence.
          Typical use $t \atoc{i}$}
\notation{\arity}{1}
         {{\rm arity}(#1)}
         {The arity of a function symbol
          Typical use $arity{f}$}


%%%%%%%%%%%%%%%% les symboles de relation

\notation{\Cc}{0}{~\|~}
         {Binds a term and a constraint to form a constraint term.
          Typical use: $(t \Cc c)$}

\notation{\EqMod}{1}{=_{#1}}
         {Equality modulo $A$ on terms}

\notation{\hEmb}{0}{\unrhd_{emb}}
         {homeomorphic embedding}
\notation{\gEmb}{1}{\unrhd_{#1}}
         {general embedding}

\notation{\equivQA}{1}{\asymp_{#1}}
         {The equivalence generated by a quasi-ordering}

\notation{\Ecl}{1}{\langle #1 \rangle_A}
         {The equivalence class modulo $A$ of a term. 
          Typical use $\Ecl{t}$.}

\notation{\Uequiv}{1}{\Lra_{#1}}
         {The equivalence between unification problems. Typical use
         $P\Uequiv{A}P'$.}

\notation{\poleqE}{0}{\equiv_E}
         {The equivalence generated by Ideal(E) on polynomials}

\notation{\ord}{0}{>}
         {ordering on (ground) terms}
\notation{\ordeq}{0}{\geq}
         {non-strict ordering on (ground) terms}
\notation{\ordE}{0}{>_{\calE}}
         {ordering on (ground) equalities}
\notation{\ordC}{0}{>_{\calC}}
         {ordering on (ground) clauses}

\notation{\subtermSt}{0}{\lhd_{sub}}
         {strict subterm relation}
\notation{\subtermEq}{0}{\unlhd_{sub}}
         {subterm relation}
\notation{\subtermStrev}{0}{\rhd_{sub}}
         {strict subterm relation reversed}
\notation{\subtermEqrev}{0}{\unrhd_{sub}}
         {subterm relation reversed}
\notation{\match}{0}{\leq}
         {subsumption ordering on terms or substitutions}
\notation{\ismatched}{0}{\geq}
         {subsumption ordering on terms or substitutions}
\notation{\smatch}{0}{<}
         {strict subsumption ordering on terms or substitutions}
\notation{\matchEquiv}{0}{\equiv}
         {subsumption equivalence of terms or substitutions}

\notation{\matchMod}{1}{\leq_{#1}}
         {subsumption ordering modulo $A$ on terms or substitutions}
\notation{\smatchMod}{1}{<_{#1}}
         {strict subsumption ordering modulo $A$ on terms or substitutions}
\notation{\matchEquivMod}{1}{\equiv_{#1}}
         {subsumption equivalence modulo on terms or substitutions}

%%%% restriction a un ensemble de variables

\notation{\matchOn}{1}{\leq^{#1}}
         {subsumption ordering for substitutions on a set of variables $V$}
\notation{\equalOn}{1}{=^{#1}}
         {equality for substitutions on a set of variables $V$}
\notation{\smatchOn}{1}{<^{#1}}
         {strict subsumption ordering for substitutions on a set of 
          variables $V$}
\notation{\matchEquivOn}{1}{\equiv^{#1}}
         {subsumption equivalence for substitutions on a set of variables $V$ }

\notation{\matchModOn}{2}{\leq^{#1}_{#2}}
         {subsumption ordering modulo a set of axioms $A$ for
          substitutions on a set of variables $V$.
          Typical use: $\matchModOn{V}{E}$ }
\notation{\equalModOn}{2}{=^{#1}_{#2}}
         {equality modulo $A$ for substitutions on a set of variables $V$}
\notation{\smatchModOn}{2}{<^{#1}_{#2}}
         {strict subsumption ordering modulo $A$ for substitutions on a set of 
          variables $V$}
\notation{\matchEquivModOn}{2}{\equiv^{#1}_{#2}}
         {subsumption equivalence modulo $A$, for substitutions on a
          set of variables $V$} 

\notation{\encompass}{0}{\sqsubseteq}
         {encompassment ordering on terms or substitutions}
\notation{\sencompass}{0}{\sqsubset}
         {strict encompassment ordering on terms or substitutions}

\notation{\encompassrev}{0}{\sqsupseteq}
         {encompassment ordering on terms or substitutions reversed}
\notation{\sencompassrev}{0}{\sqsupset}
         {strict encompassment ordering on terms or substitutions reversed}

\notation{\occurCheck}{0}{\prec^{oc}}
         {Occur check relation between variables}

%%%%%%%%%%%%%%%% les notations d'objets

\notation{\f}{2}{f(#1_1,\ldots,#1_{#2})}
         {term with $n$ subterms. Typical use: $\f{x}{n}$.}
\notation{\g}{2}{g(#1_1,\ldots,#1_{#2})}
         {term with $n$ subterms. Typical use: $\g{x}{n}$.}
\notation{\xun}{0}{x_1,\ldots,x_n}
         {sequence of $n$ variables}
\notation{\tun}{0}{t_1,\ldots,t_n}
         {sequence of $n$ terms}
\notation{\vrai}{0}{{\sf T}}
          {the symbol true}
\notation{\faux}{0}{{\sf F}}
          {the symbol false}
\notation{\Th}{1}{{\cal TH}(#1)}
          {the equational theory generated by a set of axioms $E$.
           Typical use: $\Th{E}$.}
\notation{\Mod}{1}{{\calM{}od}(#1)}
          {the class of models of a set of axioms $E$.
           Typical use: $\Mod{E}$.}

%%%%%%%%%%%%%%%% les reecritures et leur associe'es

\notation{\ra}{0}{\rightarrow}
         {rule definition symbol. Typical use $lhs \ra rhs$}
\notation{\Equiv}{1}{\stackrel{#1}{\longleftrightarrow}}
         {symetric closure of a rewriting relation. Typical use: $\equiv{*}$}
\notation{\EquivR}{2}{\stackrel{#1}{\longleftrightarrow}_{#2}}
         {symetric closure of a rewriting relation, specify moreover
         the set of axioms of rules which is used. 
         Typical use: $\equiv{*}{R}$} 
\notation{\EquivROcc}{3}{\stackrel{#1}{\longleftrightarrow}_{#2}^{#3}}
         {symetric closure of a rewriting relation, specify moreover
         the set of axioms of rules which is used and the occurrence
         of application.
         Typical use: $\equiv{*}{R}{occ}$} 
\notation{\rew}{3}{\mathop{\stackrel{#1}{\longrightarrow}^{#3}\!\!\!\!{~_{#2}}}}
         {rewriting. Typical use:  $\rew{*}{RE}{[3]}$}
\notation{\wer}{3}{\mathop{{}_{#2}^{#3}\!\!\stackrel{#1}{\longleftarrow}}}
         {reverse rewriting. Typical use: $\wer{*}{RE}{[3]}$}
\notation{\nf}{1}{#1\!\!\downarrow}
         {Normal form of a term, when it exists.
          Typical use: $nf{t}$}

%%%%%%%%%%%%%%%% forme de preuve pour la syntaxicite'

\notation{\EquivNoTop}{2}
         {\:{\stackrel{#1}{\longleftrightarrow}}_{#2}^{\neq\topoc}\:}
         {A proof below the root.
          Typical use: $t \EquivNoTop{*}{E} t'$}

\notation{\EquivAtTop}{2}
         {\:{\stackrel{#1}{\longleftrightarrow}}_{#2}^{\topoc}\:}
         {A proof step at the top occurrence.
          Typical use: $t \EquivAtTop{\delta}{E} t'$}

\notation{\EquivOneTop}{3}
         {\:{\stackrel{#1}{\longleftrightarrow}}_{#2}^{#3\topoc}\:}
         {A proof with at most one application at $\topoc$.
          Typical use: $t \EquivOneTop{*}{E}{\delta} t'$}

\notation{\dac}{1}
         {\downarrow_{#1}^{\topoc}}
         {Typical use: $l = f(\vec r) \dac{E} f(\vec g) = d$}

\notation{\soul}{1}
         {\underline{#1}}
         {Underline in order to constantify: Typical use $\soul{t}$}

%%%%%%%%%%%%%%%% La surreduction

\notation{\surred}{2}{\leadsto^{#1}_{#2}}
         {one narrowing step. Typical use:  $\surred{RE}{[3]}$}

\notation{\Surred}{3}{\stackrel{#1}{\leadsto^{#2}_{#3}}}
         {Several narrowing steps. Typical use:  $\Surred{*}{RE}{[3]}$}

\notation{\bsurred}{2}{~^{b}\!\!\!\leadsto^{#1}_{#2}}
         {one basic narrowing step. Typical use:  $\bsurred{RE}{[3]}$}

\notation{\csurred}{2}{~^{c}\!\!\!\leadsto^{#1}_{#2}}
         {one constraint narrowing step. Typical use:  $\csurred{RE}{[3]}$}



%%%%%%%%%%%%%%%% les regles conditionnelles et clauses de Horn

\notation{\Ci}{0}{{\msam V}}
         {The implication for conditional rewriting}

\notation{\condeq}{3}{{#1} \, \Ci \,{#2} = {#3}}
         {egalite conditionnelle}
\notation{\cond}{3}{{#1} \, \Ci \,{#2} \rightarrow {#3}}
         {regle conditionnelle}
\notation{\horn}{2}{{#1} \Ci {#2}}
         {clause de Horn}

%%%%%%%%%%%%%%%% les objets avec contraintes

\notation{\Con}{0}{\parallel}
         {The context for constraint rewriting}

\notation{\CT}{3}{(\exists {#1}, {#2} \Cc {#3})}
         {Constraint term. Typical use $\CT{X}{t}{c}$}
\notation{\consteq}{3}{{#1} = {#2} \, \Con \, {#3}}
         {egalite contrainte}
\notation{\constrr}{3}{{#1} \rightarrow {#2} \, \Con \, {#3}}
         {regle contrainte}
\notation{\constcl}{2}{{#1} \, \Con \, {#2}}
         {clause contrainte}


%%%%%%%%%%%%%%%% les operations sur les termes

\notation{\replace}{3}{#1[#2 \la #3]}
         {replacement in a term at some occurrence. 
          Typical use $\replace{t}{m}{t'}$}

%%%%%%%%%%%%%%%% les comparaisons d'occurrences
\notation{\incomp}{0}{\Join}
         {incompatible positions. Typical use $\omega_1 \incomp \omega_2$}

